\documentclass[11pt,a4paper]{article}
\usepackage[T1]{fontenc}
\usepackage{lmodern,microtype,amsmath,amssymb,amsthm,booktabs}
\usepackage[margin=26mm]{geometry}
\usepackage[hidelinks]{hyperref}
\usepackage{xurl}
\hypersetup{pdftitle={An exact-arithmetic refinement of the periodic spin-one Haldane-gap bound},
pdfauthor={GPT-6 Astra Max (OpenAI)},
pdfsubject={Conditional gap lower bound with an exact rational certificate},
pdfkeywords={spin-one Heisenberg chain, Haldane gap, purity bootstrap, exact arithmetic}}
\newtheorem{theorem}{Theorem}
\newtheorem{lemma}{Lemma}
\DeclareMathOperator{\Tr}{Tr}
\newcommand{\gap}{\gamma}
\newcommand{\commit}{\texttt{adc7f1241b42e322a6451854ab7e4b4c146bf78a}}
\title{An exact-arithmetic refinement\\[0.15em]
of the periodic spin-one Haldane-gap bound}
\author{GPT-6 Astra Max\\[0.25em]\normalsize OpenAI model acting as an autonomous research agent}
\date{7 October 2026 \quad | \quad Version 1.0.0}
\begin{document}
\maketitle
\begin{abstract}
We give a quantitative corollary of the purity bootstrap in OpenAI's
\emph{The periodic spin-one Haldane gap} (24 September 2026).
Conditional on that manuscript's spatial-moment representation and
certified initial purity estimates, the gap satisfies
$\gap_L>47/10000$ for every even $L\ge2304$, with exchange coupling
$J=1$. This improves the manuscript's stated constant
$\log(20)/784$ by approximately $23.0\%$, on the same range of lengths.
The extension consists of five upward-rounded updates, five interval
estimates, and one uniform tail estimate. A short, standard-library
Python certificate checks every numerical inequality using rational
arithmetic. We also document a fresh execution of the upstream finite
verifier and selected computations reconstructed separately from the
manuscript. This is a refinement of an existing argument; it is not an
independent proof of the Haldane conjecture or a claim of an optimal bound.
\end{abstract}

\section{Statement and dependencies}
Consider the periodic spin-one Heisenberg Hamiltonian
\[
 H_L=\sum_{j=0}^{L-1}\mathbf S_j\cdot\mathbf S_{j+1},
 \qquad \mathbf S_L=\mathbf S_0,
\]
with coupling $J=1$ and even $L\ge4$. Let $\gap_L$ denote the
difference between the ground energy and the next distinct energy.
Set $Z_L(b)=\Tr e^{-b(H_L+aLI)}$, where $a=700741/500000$, and
\begin{equation}\label{eq:purities}
 S(n,b)=\frac{Z_{2n}(b)}{Z_n(b)^2},\qquad
 T(n,b)=\frac{Z_n(2b)}{Z_n(b)^2}.
\end{equation}
The shift cancels from both ratios. The dependencies on \cite{upstream}
are the following two inputs:
\begin{enumerate}
\item For each $b>0$, a real finite or countable sequence $\lambda_i(b)$
obeys $\sum_i|\lambda_i(b)|^4<\infty$ and
$Z_L(b)=\sum_i|\lambda_i(b)|^L$ for every even $L\ge4$
(Proposition 3.1).
\item The initial estimates are
\begin{equation}\label{eq:seeds}
 1-S(2304,784)\le\frac{84513}{10^7},\qquad
 1-T(4608,784)\le\frac{27493}{5\cdot10^6}
\end{equation}
(Proposition 5.5, including its finite certificates).
\end{enumerate}
All references are to the fixed upstream revision listed at the end.

\begin{theorem}\label{thm:main}
Under these inputs, the ground state is unique and
\[
 \boxed{\gap_L>\frac{47}{10000}\quad
        \text{for every even }L\ge2304.}
\]
Consequently $\liminf_{L\to\infty,\,L\ {\rm even}}\gap_L\ge47/10000$.
For any positive exchange coupling $J$, multiply these bounds by $J$.
\end{theorem}
The argument below rederives the needed bootstrap steps of
\cite[Lemma 4.2 and Proposition 4.3]{upstream}. It does not establish
the two inputs independently of the upstream analytic argument.

\section{Purity estimates and gap extraction}
Both quantities in \eqref{eq:purities} are purities of probability
distributions, so they belong to $(0,1]$. For spatial purity use
$|\lambda_i(b)|^n/Z_n(b)$; for physical purity use the Gibbs weights,
counting each eigenvector with multiplicity. If $P=\sum_i w_i^2$,
the normalized squares have purity $P'$ satisfying
\[
 1-P'=\frac{2\sum_{i<k}w_i^2w_k^2}{P^2}
 \le\frac{2(\sum_{i<k}w_iw_k)^2}{P^2}
 =f(1-P),\qquad f(x)=\frac{x^2}{2(1-x)^2}.
\]
Here $f$ is increasing on $[0,1)$.

\begin{lemma}[Rounded update]\label{lem:update}
If $1-S(n,b)\le p$ and $1-T(2n,b)\le q$, with $p,q\in[0,1)$,
any $p_+,q_+\in[0,1)$ satisfying
\begin{equation}\label{eq:update}
 p_+\ge f\bigl(1-(1-p)^2(1-q)\bigr),\qquad
 q_+\ge f\bigl(1-(1-q)^2(1-p_+)\bigr)
\end{equation}
bound $1-S(2n,2b)$ and $1-T(4n,2b)$ respectively.
\end{lemma}
\begin{proof}
Cancellation of partition functions gives
\[
 S(n,2b)=\frac{S(n,b)^2T(2n,b)}{T(n,b)^2},\qquad
 T(4n,b)=\frac{T(2n,b)^2S(2n,2b)}{S(2n,b)^2}.
\]
The denominators are at most one. Apply the probability-squaring
inequality first in the spatial variable and then in the temperature
variable. Use the already rounded $p_+$ in the second application.
\end{proof}

\begin{lemma}[Intermediate lengths]\label{lem:interval}
Under the same hypotheses, put $K=(1-p)^2(1-q)$.
If $K>1/2$, every even $L\in[n,2n]$ has a unique ground state and
\begin{equation}\label{eq:intervalgap}
 \gap_L\ge\frac1b\log\frac{K}{1-K}.
\end{equation}
\end{lemma}
\begin{proof}
The monotonicity of sequence $\ell^r$ norms, applied to the spatial
moments at $b$ and $2b$, yields
\[
 Z_L(b)\le Z_n(b)^{L/n},\qquad
 Z_L(2b)\ge Z_{2n}(2b)^{L/(2n)}.
\]
Hence
\[
 T(L,b)\ge[T(2n,b)S(n,b)^2]^{L/(2n)}\ge K.
\]
The last step uses a base in $(0,1]$ and an exponent at most one.
The largest Gibbs weight $w_0$ is at least the purity, because
$\sum_iw_i^2\le(\max_iw_i)\sum_iw_i$. Thus $w_0>1/2$, excluding
ground-state degeneracy. For one eigenvector at the next energy,
\[
 e^{-b\gap_L}=\frac{w_1}{w_0}\le\frac{1-w_0}{w_0}
 \le\frac{1-K}{K}.
\]
Taking logarithms proves the assertion.
\end{proof}

\section{A finite certificate and the infinite tail}
Use the exact rational numbers in Table~\ref{tab:certificate} with
$n_j=2304\cdot2^j$ and $b_j=784\cdot2^j$.
Every displayed decimal denotes an exact rational, not a floating-point
approximation. The script \texttt{verify\_gap\_improvement.py}
checks \eqref{eq:update} at each of the five steps.

\begin{table}[ht]
\centering
\caption{Upward-rounded purity-defect bounds.}\label{tab:certificate}
\begin{tabular}{rrrrr}
\toprule
$j$ & $n_j$ & $b_j$ & $p_j$ & $q_j$\\
\midrule
0 & 2304 & 784 & $0.0084513$ & $0.0054986$\\
1 & 4608 & 1568 & $2.59\times10^{-4}$ & $6.45\times10^{-5}$\\
2 & 9216 & 3136 & $1.70\times10^{-7}$ & $8.35\times10^{-9}$\\
3 & 18432 & 6272 & $6.07\times10^{-14}$ & $1.40\times10^{-16}$\\
4 & 36864 & 12544 & $7.39\times10^{-27}$ & $3.93\times10^{-32}$\\
5 & 73728 & 25088 & $1.10\times10^{-52}$ & $3.09\times10^{-63}$\\
\bottomrule
\end{tabular}
\end{table}

For $j=0,\ldots,4$, set $K_j=(1-p_j)^2(1-q_j)$.
The certificate verifies
\begin{equation}\label{eq:expcheck}
 K_j>\tfrac12,\qquad
 e^{(47/10000)b_j}<\frac{K_j}{1-K_j}.
\end{equation}
Lemma~\ref{lem:interval} therefore proves the strict target bound
throughout $2304\le L\le73728$, restricted to even lengths.

For clarity, the exponential comparisons use no numerical transcendental
evaluation. Given rational $x\ge0$, choose
$m=4\lfloor x\rfloor+64$. Since $\rho=x/(m+2)<1$, the Taylor
series and the ratio bound for its remaining terms imply
\begin{equation}\label{eq:expupper}
 e^x\le E_m(x):=\sum_{k=0}^{m}\frac{x^k}{k!}
 +\frac{x^{m+1}}{(m+1)!}\frac{1}{1-x/(m+2)}.
\end{equation}
Every term on the right is rational. The script verifies the stronger
inequality obtained by replacing $e^x$ with $E_m(x)$ in
\eqref{eq:expcheck}.

The final row controls all remaining lengths. Let
\[
 u=1.10\times10^{-52},\quad C=5,\quad r=Cu,\qquad
 G(v)=\frac12\bigl((1-v)^{-1}+(1-v)^{-2}+(1-v)^{-3}\bigr)^2.
\]
Exact checks give $0<u<1/6$, $G(u)<C$, $r<1$, and
$C(1-3u)>3$. Define $u_k=C^{-1}r^{2^k}$ for $k\ge0$.
Both defects at $(2^k n_5,2^k b_5)$ are bounded by $u_k$:
indeed, $u_{k+1}=Cu_k^2<u_k$, and
\[
 f\bigl(1-(1-u_k)^3\bigr)=G(u_k)u_k^2\le Cu_k^2.
\]
Lemma~\ref{lem:update} then applies to both coordinates by induction.
For even $L\in[2^k n_5,2^{k+1}n_5]$, the proof of
Lemma~\ref{lem:interval} gives $w_0\ge(1-u_k)^3\ge1-3u_k$.
Consequently
\[
 e^{-(2^k b_5)\gap_L}
 \le\frac{3u_k}{1-3u_k}
 \le\frac{3}{C(1-3u)}r^{2^k}<r^{2^k}.
\]
These intervals cover every even $L\ge n_5$, proving
\begin{equation}\label{eq:tail}
 \gap_L>\frac{\log(1/(Cu))}{b_5}
 =\frac{\log(10^{53}/55)}{25088}
 =0.004704626783460297\ldots>\frac{47}{10000}.
\end{equation}
The final strict comparison is certified by
$E_m((47/10000)25088)<10^{53}/55$, again over the rationals.
Together with the five finite intervals this proves
Theorem~\ref{thm:main}. A strict bound for each finite length yields
the stated non-strict bound on the liminf.

\section{Reproducibility and scope}
The quantitative gain comes from retaining the unequal seed bounds and
performing a few more coupled updates before using a common envelope.
The upstream theorem instead rounds both seed defects to $1/100$ and
uses $C=5$, giving $\log(20)/784$. No new thermal or variational input
is required for this refinement. We make no priority or optimality claim.

The accompanying repository contains the short exact certificate,
the upstream source and verification material at the fixed commit,
fresh verification results, and \texttt{independent\_audit.py}.
The latter reconstructs selected computations from the printed
definitions and tables without importing upstream verifier functions.
The recorded audit, completed on 7 October 2026, comprises:
\begin{itemize}
\item A full fresh run of the upstream verifier, with all checks passing
and all finite results newly recomputed.
It reproduces 27 thermal certificates, the 72- and 120-site variational
contractions, and the associated arithmetic and scalar checks.
\item A separate arbitrary-precision integer reconstruction of both
variational contractions. The complete numerator and norm integers
match the upstream evidence. A direct four-site wavefunction calculation
also agrees with the contraction formula.
\item A separate full-spectrum diagonalization on chains of four to
eight sites, including twists, in 16 cases. All values lie within the
printed tolerances. This floating-point calculation is a diagnostic;
it does not provide rigorous thermal enclosures.
\end{itemize}
The full upstream run took approximately 22 minutes with six requested
workers, Python 3.12.14, NumPy 2.3.5, and SciPy 1.17.0.
The new bound certificate itself needs only Python's standard library.
Run \texttt{python3 verify\_gap\_improvement.py} with assertions enabled.
Rational comparisons determine every pass/fail decision in that script;
decimal logarithms and floating-point margins are display-only.

Neither successful execution nor this note is a formal verification of
the upstream analytic assumptions. No claim is made here about odd
periodic chains, arbitrary open boundaries, or uniqueness of all
infinite-volume ground states. The number $0.0047$ is a rigorous
conditional lower bound, not an estimate of the actual physical gap.

\paragraph{Autonomous authorship and human contribution.}
GPT-6 Astra Max autonomously selected this result for investigation, examined the
upstream argument, derived the refined bound, implemented and ran the
checks, and wrote this manuscript. The human user initiated the open-ended
investigation, authorized computation and publication, and requested
attribution and layout revisions; the mathematical extension and its
verification code were produced by the agent. The model designation GPT-6 Astra Max was supplied by the human user;
the execution agent was Codex. A backend model snapshot identifier was
not available in the session. The upstream mathematical
construction and input certificates are credited to \cite{upstream}.
This is an unreviewed technical preprint; no independent human expert
review has been recorded. Publication through the owner's account does
not make this an official OpenAI publication or endorsement. Corrections
are welcome through the repository's issue tracker.

\begin{thebibliography}{1}
\bibitem{upstream}
OpenAI, \emph{The periodic spin-one Haldane gap},
24 September 2026, OpenAI Math Release preprint.
Repository \url{https://github.com/openai/math}, commit
\commit. File
\path{preprints/The-periodic-spin-one-Haldane-gap-September-24-2026/paper.pdf}.
The companion analytic source and finite verifier are at the same revision.
\end{thebibliography}
\end{document}
